(1) 记
\[
t_c=\tanh\frac{\hbar\omega_c}{2k_BT_c},\qquad
t_h=\tanh\frac{\hbar\omega_h}{2k_BT_h},\qquad
a=1-2\xi.
\]
在能量本征基上 \(z=\langle\sigma_z\rangle=-t\)，对角能量平均为 \(\hbar\omega z/2\)。有限时间孤立冲程的两能级跃迁矩阵为
\(\begin{pmatrix}1-\xi&\xi\\\xi&1-\xi\end{pmatrix}\)，故对角极化变为 \(z\mapsto az\)，因此
\[
U_A=-\frac{\hbar\omega_c}{2}t_c,\quad
U_B=-\frac{\hbar\omega_h}{2}a t_c,\quad
U_C=-\frac{\hbar\omega_h}{2}t_h,\quad
U_D=-\frac{\hbar\omega_c}{2}a t_h.
\]
热等容冲程没有频隙变化造成的机械功，故
\[
Q_h=U_C-U_B=\frac{\hbar\omega_h}{2}(a t_c-t_h),\quad
Q_c=U_A-U_D=\frac{\hbar\omega_c}{2}(a t_h-t_c).
\]
\(Q_h>0\) 是作为热机工作的必要条件。

(2) 一个完整周期 \(\Delta U=0\)，所以
\[
W_{\rm out}=Q_h+Q_c
=\frac\hbar2\big[\omega_h(a t_c-t_h)+
\omega_c(a t_h-t_c)\big].
\]
正功热机需 \(Q_h>0\) 且 \(W_{\rm out}>0\)，其中第二式可直接代数检验；在本设定 \(0\le a\le1\) 与 \(T_h>T_c,\omega_h>\omega_c\) 下，正功条件实际自动意味着 \(Q_h>0\)。
效率为
\[
\eta=\frac{W_{\rm out}}{Q_h}
=1-\frac{\omega_c(t_c-a t_h)}{\omega_h(a t_c-t_h)}.
\]
无跃迁时 \(a=1\)，只要 \(t_c>t_h\)，就有
\(\eta_0=1-\omega_c/\omega_h\)。
比较相同两热浴及频隙但 \(\xi>0\)，得到严格损功
\[
W_{\rm out}(0)-W_{\rm out}(\xi)
=\hbar\xi(\omega_h t_c+\omega_c t_h)>0.
\]
它来自孤立冲程内部的非绝热布居跃迁，而非等容冲程对外机械做功。

(3) 熵产生的量纲是能量/温度，循环工作物质熵返回初始值，故
\[
\Sigma=-\frac{Q_h}{T_h}-\frac{Q_c}{T_c}.
\]
设 \(\rho_B,\rho_D\) 为压缩、膨胀后的态，\(\rho_C,\rho_A\) 为对应热态。每个孤立冲程是幺正演化，因此 \(S_{\rm vN}(\rho_B)=S_{\rm vN}(\rho_A)\)，
\(S_{\rm vN}(\rho_D)=S_{\rm vN}(\rho_C)\)。两热化过程相对熵恒满足
\[
k_B D(\rho_B\Vert\rho_C)
+k_B D(\rho_D\Vert\rho_A)
=-\frac{Q_h}{T_h}-\frac{Q_c}{T_c}=\Sigma\ge0.
\]
从 \(\Sigma\ge0\) 且 \(Q_h>0\) 得
\(Q_c\le-(T_c/T_h)Q_h\)，故
\(\eta\le1-T_c/T_h\)。
对 \(\xi=0\)，\(Q_h>0\iff t_c>t_h\iff
\omega_c/T_c>\omega_h/T_h\iff
\omega_c/\omega_h>T_c/T_h\)，所以正功Otto效率确实严格小于Carnot界。

(4) 从第二问可见
\(dW_{\rm out}/d\xi=-\hbar(\omega_h t_c+\omega_c t_h)<0\)。
当无跃迁理想净功 \(W_0=\hbar(\omega_h-\omega_c)(t_c-t_h)/2>0\) 时，临界跃迁概率为
\[
\xi_*=\frac{(\omega_h-\omega_c)(t_c-t_h)}
{2(\omega_h t_c+\omega_c t_h)}.
\]
\(0\le\xi<\xi_*\) 为正功区，\(\xi=\xi_*\) 净功为零，之后进入需外部供功状态。令
\(A=t_c-t_h>0\)，小 \(\xi\) 时
\[
\eta=\left(1-\frac{\omega_c}{\omega_h}\right)
-\frac{2\omega_c}{\omega_h}
\frac{\xi(t_c+t_h)}{A}+O(\xi^2),
\]
并非简单仍等于无摩擦Otto效率。作为能量核验，可取无量纲
\(\hbar=k_B=1\)、\((\omega_c,T_c)=(1,1)\)、\((\omega_h,T_h)=(2,3)\)、
\(\xi=0.01\)，从第一问逐项计算后必得
\(W_{\rm out}=Q_h+Q_c\) 且相对 \(\xi=0\) 输出功下降
\(0.01[2\tanh(1/2)+\tanh(1/3)]>0\)。
